\section{The generic quintic has Galois group \texorpdfstring{$S_5$}{S5}}
\label{sec:generic-quintic}

The usual proof is often presented through the theorem on symmetric rational
functions.  That theorem is not needed here.  We define the coefficient field
directly as the subfield generated by the elementary symmetric functions and
identify the automorphism group by its action on the five roots.

\subsection{The fields and the polynomial}

Let $x_1,\ldots,x_5$ be algebraically independent indeterminates over $\Q$ and
put
\[
 L=\Q(x_1,x_2,x_3,x_4,x_5).
\]
For $1\le j\le5$, let
\[
 e_j=\sum_{1\le i_1<\cdots<i_j\le5}
       x_{i_1}\cdots x_{i_j}
\]
be the elementary symmetric functions, and let
\[
 K=\Q(e_1,e_2,e_3,e_4,e_5)\subseteq L.
\]
Define
\[
 P(T)=\prod_{i=1}^{5}(T-x_i).
 \tag{G1}\label{eq:generic-product}
\]
Expanding by elementary symmetric functions gives
\[
 P(T)=T^5-e_1T^4+e_2T^3-e_3T^2+e_4T-e_5\in K[T].
 \tag{G2}\label{eq:generic-coefficients}
\]

The elements $x_i$ are pairwise distinct.  Indeed, $x_i-x_j$ is a nonzero
polynomial when $i\ne j$, hence a nonzero element of the fraction field $L$.
Therefore $P$ is separable.  Each $x_i$ is algebraic over $K$ because it
satisfies the monic polynomial $P$, so the extension
$K(x_1,\ldots,x_5)/K$ is finite.  The polynomial splits in $L$, and
\[
 L=K(x_1,\ldots,x_5),
\]
so $L$ is the finite splitting field of $P$ over $K$.  Hence $L/K$ is normal
and separable, and its $K$-automorphism group is its Galois group.

All these claims have direct sparse-polynomial certificates.  Algebraic
independence is built into the representation of $L$ as the fraction field of
the polynomial ring $\Q[x_1,\ldots,x_5]$; no existence theorem for a generic
point is required.

\subsection{Every permutation gives an automorphism}

For $\pi\in S_5$, define a $\Q$-automorphism $\widehat\pi$ of $L$ by
\[
 \widehat\pi(x_i)=x_{\pi(i)}.
\]
On a rational function, apply the same variable renaming to its numerator and
denominator.  This is a primitive-recursive operation on sparse codes.
Because every $e_j$ is symmetric,
\[
 \widehat\pi(e_j)=e_j,
\]
so $\widehat\pi$ fixes $K$ pointwise.  Hence
\[
 S_5\longrightarrow\Gal(L/K),
 \qquad
 \pi\longmapsto\widehat\pi
 \tag{G3}\label{eq:S5-to-Gal}
\]
is a homomorphism.  It is injective: if $\pi(i)\ne\rho(i)$, then
$x_{\pi(i)}\ne x_{\rho(i)}$, so the two automorphisms differ on $x_i$.

\subsection{Every automorphism is a permutation of the roots}

Let $\sigma\in\Gal(L/K)$.  Since $\sigma$ fixes the coefficients of $P$,
\[
 P(\sigma(x_i))=\sigma(P(x_i))=0.
 \tag{G4}\label{eq:auto-root}
\]
But in $L$ we have the exact factorization~\eqref{eq:generic-product}.  Hence
\[
 \prod_{j=1}^{5}(\sigma(x_i)-x_j)=0.
\]
A field has no zero divisors, so $\sigma(x_i)=x_j$ for some $j$.  The root is
unique because the $x_j$ are pairwise distinct.  Thus there is a function
$\pi_\sigma:\{1,\ldots,5\}\to\{1,\ldots,5\}$ with
\[
 \sigma(x_i)=x_{\pi_\sigma(i)}.
\]
The function is injective because $\sigma$ is injective; a bounded finite
count proves that an injection on five elements is a bijection.  Therefore
$\pi_\sigma\in S_5$.

The automorphisms $\sigma$ and $\widehat{\pi_\sigma}$ agree on $\Q$ and on all
five generators $x_i$ of $L$, so they are equal.  Consequently
\eqref{eq:S5-to-Gal} is surjective as well as injective.

\begin{theorem}[Generic Galois group]
\label{thm:generic-galois-S5}
For the polynomial $P$ in~\eqref{eq:generic-coefficients},
\[
 \Gal(L/K)\cong S_5.
\]
The isomorphism and its inverse are given explicitly by variable renaming and
by reading the images of the five generators.
\end{theorem}

\begin{proof}
The preceding two subsections construct mutually inverse homomorphisms.  Every
step is an equality of rational functions or a bounded calculation on five
indices.
\end{proof}

\subsection{Why no fixed-field theorem is being hidden}

One may instead define $K$ as the fixed field $L^{S_5}$ and use Artin's
fixed-field theorem.  ProveIt's existing \code{GenericAbelRuffini.lean} follows
that elegant route.  For the iPRA arithmetization, the explicit field
$\Q(e_1,\ldots,e_5)$ is convenient because Theorem
\ref{thm:generic-galois-S5} needs only:

\begin{itemize}
\item sparse polynomial expansion of a five-factor product;
\item primitive-recursive renaming of five variables;
\item exact rational-function equality;
\item the no-zero-divisors rule; and
\item finite injection-to-bijection on a five-element list.
\end{itemize}

The fundamental theorem of symmetric polynomials, linear independence of
field embeddings, and degree counting are all unnecessary for this
identification.

\subsection{The same construction in every degree}

For any $n\ge2$, the identical argument with
$L_n=\Q(x_1,\ldots,x_n)$ and
$K_n=\Q(e_1,\ldots,e_n)$ gives
\[
 \Gal\left(L_n/K_n\right)\cong S_n
\]
for the generic degree-$n$ product.  The proof below only needs $n=5$ because
$S_5$ is already nonsolvable.  Higher degrees can therefore be handled either
by the generic $S_n$ construction or, more economically, by padding the
quintic.
